\section{问题重述}

\subsection{问题背景}

研究对象是一个由小区负载、光伏发电和储能设备组成的微网。微网既可以利用本地光伏，也可以从外部电网购电；储能则把不同时刻的电能联系起来，使系统能够把低价时段或光伏富余时段的电量转移到后续使用。由于居民负载、光伏出力以及外网价格都随时间变化，购电决策不能只考虑某一个时刻，而要同时满足当前供电、后续储能状态和全天费用三个方面的要求。

题目以 $10$ min 为基本时间尺度。储能额定容量为 $12000$ kWh，运行电量必须保持在 $1200\sim10800$ kWh 之间，最大充、放电功率均为 $5000$ kW，充放电效率均按 $90\%$ 处理。因而，一天的 $144$ 个调度时段通过储能状态递推相互耦合：某一时段多充入的电，会改变后续可放出的电量；某一时段过度放电，也可能使之后的高价时段失去调节能力。

除物理约束外，本题还有一个贯穿四问的关键：\textbf{每次决策只能使用该时刻已经获得的信息。} 问题一的信息基本确定，问题二需要面对全天计划与实际运行之间的偏差，问题三进一步允许在日内根据新预报修改尚未执行的计划，问题四则把固定典型日电价推广为逐日波动电价。因此，四问可看作同一微网调度问题在不同信息条件和结算规则下的逐步扩展。

\subsection{问题提出}

\indent\textbf{问题一：}在典型日电价、负载和光伏预测均给定的条件下，于每日 $0{:}00$ 制定全天计划购电与储能充放电策略，使每个时段均能满足小区负荷，并保证日初与日末储电量相同。在论文中给出题目指定时段的购电量、全天购电量与购电费，以及各 $4$ h 时段的储能充放电量；完整方案写入 \texttt{result1.xlsx}。

\indent\textbf{问题二：}当负载和光伏随日期变化时，每天仍需在 $0{:}00$ 先制定全天计划。若实际运行中供电不足，缺口需按交易时刻正常电价的 $5$ 倍紧急购入。因此需要研究：在未来源荷尚未实现、预测存在误差时，如何利用历史数据制定既不过度冒险、又不过度保守的购电策略，并用储能吸收日内偏差。按题目要求给出四个指定日期的购电、储能与紧急购电结果，完整方案写入 \texttt{result2.xlsx}。

\indent\textbf{问题三：}每天 $0{:}00$、$6{:}00$、$12{:}00$、$18{:}00$ 均可获得未来 $24$ h 的光伏预报。$0{:}00$ 制定全天计划后，后续节点可修改尚未执行的购电量；向上调整和向下调整采用不同的结算规则。因此需要建立滚动优化模型，并通过对照实验回答“是否值得引入其他时刻的预报”。完整结果写入 \texttt{result3.xlsx}。

\indent\textbf{问题四：}将问题二、三推广到波动电价。根据附件4给出的逐日实时电价重新计算两类策略，比较固定典型日电价与波动电价下的购电、储能和紧急购电结果，完整方案分别写入 \texttt{result4-2.xlsx} 与 \texttt{result4-3.xlsx}。
